% !TEX TS-program = lualatex

\documentclass[border = 3pt]{standalone}

%%%
% Des couleurs nommées faciles d'emploi via le package ''xcolor''!
%%%
\usepackage[svgnames]{xcolor}

%%%
% La bibliothèque \luadraw allie une facilité d’utilisation
% à un rendu particulièrement soigné.
%%%
\usepackage[3d]{luadraw}


%%%
% La bibliothèque ''shadings'' utilisable via \tikz est nécessaire
% pour faire fonctionner ''luadraw_shadedforms''. .
%%%
\usetikzlibrary{shadings}

\begin{document}

\begin{luadraw}{name = vector-field-two-grids}
require 'luadraw_palettes'
require 'luadraw_shadedforms'
require 'luadraw_fields'

------
-- ¨Def du ¨chp.
------
local f = function(x, y)
  return {1, x^2 + y^2 - 1}
end

------
-- ¨Def du module du ¨chp.
------
local f_modulus = function(x, y)
  local A = f(x, y)

  return cpx.abs(Z(A[1],A[2]))
end

------
-- Réglages liés au tracé.
------
local xmin, xmax, ymin, ymax = -3, 3, -3, 3

------
-- ¨Def de la zone graphique.
------
local graphview = graph:new{
  window = {xmin - 0.5, xmax + 1.5, ymin - 0.5, ymax + 0.5},
  size   = {10, 10},
  bbox   = false
}

------
-- Dessinons !
------
graphview:Dshadedrectangle(
  xmin, xmax, ymin, ymax,
  palJet,
  {
    values = f_modulus,
    bar    = "right",
    baroptions
  }
)

graphview:Dvectorfield(
  f,
  {
    view         = {xmin, xmax, ymin, ymax},
    draw_options = "-stealth"
  }
)

graphview:Dgradbox(
  {
    Z(xmin, ymin),
    Z(xmax, ymax),
    1, 1
  },
  {
    title = "Vector field $f(x,y) = (1, x^2 + y^2 - 1)$"
  }
)

------
-- Montrons le résultat de notre oeuvre.
------
graphview:Show()
\end{luadraw}

\end{document}
